\documentclass{article}
\usepackage{amsmath,amssymb}
\usepackage[paperwidth=440bp,paperheight=340bp,margin=24bp]{geometry}
\pagestyle{empty}
\setlength{\parindent}{0pt}
\newcommand*\pdfcraftbegin[1]{\special{pdfcraft:eq:#1:begin}}
\newcommand*\pdfcraftend[1]{\special{pdfcraft:eq:#1:end}}
\newcommand*\pdfcrafttag[1]{\special{pdfcraft:tag:#1}}
% amsmath's \maketag@@@, with the equation number bracketed.
\makeatletter
\def\maketag@@@#1{\hbox{\m@th\normalfont\pdfcrafttag{begin}#1\pdfcrafttag{end}}}
\makeatother
\relpenalty=10000
\binoppenalty=10000
\sloppy
\begin{document}
\onecolumn
Consider a point on a circle, so that $\pdfcraftbegin{sup_squares@onecol_a}
x^2 + y^2 = r^2
\pdfcraftend{sup_squares@onecol_a}$ holds for its coordinates. Halving and
thirding a cake gives $\pdfcraftbegin{frac_text_style@onecol_a}
\frac{1}{2} + \frac{1}{3} = \frac{5}{6}
\pdfcraftend{frac_text_style@onecol_a}$ of it, which is easy to check. A geometric series
$\pdfcraftbegin{sum_inline@onecol_a}
\sum_{k=0}^{\infty} z^k
\pdfcraftend{sum_inline@onecol_a}$ converges when the ratio is small, and the most famous relation of all reads
\begin{equation}\pdfcraftbegin{equation_numbered@onecol_a}
E = mc^2
\pdfcraftend{equation_numbered@onecol_a}\end{equation}
where the constant is the speed of light. Nothing more is needed here, but the reader may wish
to compare the last display with the one on the next page.

\clearpage
\onecolumn
The integral $\pdfcraftbegin{int_inline@onecol_b}
\int_a^b g(t)\,dt
\pdfcraftend{int_inline@onecol_b}$ measures the area under the curve, and the limit
$\pdfcraftbegin{lim_inline@onecol_b}
\lim_{x \to 0} \frac{\sin x}{x} = 1
\pdfcraftend{lim_inline@onecol_b}$ is the first one met in a course on calculus. When two points are close we write
$\pdfcraftbegin{abs_epsilon@onecol_b}
|x - y| < \varepsilon
\pdfcraftend{abs_epsilon@onecol_b}$ and move on. The Kronecker symbol is defined by
\begin{equation}\pdfcraftbegin{cases_numbered@onecol_b}
\delta_{ij} = \begin{cases} 1 & i = j \\ 0 & i \ne j \end{cases}
\pdfcraftend{cases_numbered@onecol_b}\end{equation}
and it appears in almost every computation with coordinates. A short aligned computation
closes the page:
\begin{align}\pdfcraftbegin{align_two@onecol_b}
a &= b + c \\ d &= e - f
\pdfcraftend{align_two@onecol_b}\end{align}

\clearpage
\twocolumn
The squaring map $\pdfcraftbegin{mapsto@twocol_a}
x \mapsto x^2
\pdfcraftend{mapsto@twocol_a}$ is a first example of a map $\pdfcraftbegin{arrow_map@twocol_a}
f\colon X \to Y
\pdfcraftend{arrow_map@twocol_a}$ between sets, which
sends each element somewhere. Greek letters such as $\pdfcraftbegin{greek_lower@twocol_a}
\alpha + \beta + \gamma + \delta
\pdfcraftend{greek_lower@twocol_a}$ name angles in most books. The integers sit
in a chain of number systems, and the next display is numbered:
\begin{equation}\pdfcraftbegin{equation_numbered@twocol_a}
E = mc^2
\pdfcraftend{equation_numbered@twocol_a}\end{equation}
which ends the first column.
\newpage

In the second column a list $\pdfcraftbegin{ldots_list@twocol_a}
x_1, \ldots, x_n
\pdfcraftend{ldots_list@twocol_a}$ of unknowns is solved by elimination, and the
relations $\pdfcraftbegin{rel_order@twocol_a}
a \le b \ge c \ne d
\pdfcraftend{rel_order@twocol_a}$ are read from left to right. A sum $\pdfcraftbegin{cdots_sum@twocol_a}
1 + 2 + \cdots + n
\pdfcraftend{cdots_sum@twocol_a}$ of the first numbers
was found by a schoolboy.

\clearpage
\twocolumn
Logarithms turn products into sums: $\pdfcraftbegin{log_product@twocol_b}
\log(xy) = \log x + \log y
\pdfcraftend{log_product@twocol_b}$ for positive numbers. The determinant
$\pdfcraftbegin{det_char@twocol_b}
\det(\lambda I - A)
\pdfcraftend{det_char@twocol_b}$ gives the characteristic polynomial, and the binomial coefficient $\pdfcraftbegin{binom_inline@twocol_b}
\binom{2m}{m}
\pdfcraftend{binom_inline@twocol_b}$
counts balanced paths. Two equations, numbered, follow:
\begin{gather}\pdfcraftbegin{gather_rows@twocol_b}
a = b \\ c = d + e
\pdfcraftend{gather_rows@twocol_b}\end{gather}
\newpage

Square roots such as $\pdfcraftbegin{sqrt_inline@twocol_b}
\sqrt{n}\log n
\pdfcraftend{sqrt_inline@twocol_b}$ appear in running times. The expectation $\pdfcraftbegin{expectation@twocol_b}
\mathbb{E}[X] = \mu
\pdfcraftend{expectation@twocol_b}$
of a random variable is its mean, and a ratio $\pdfcraftbegin{slash_fraction@twocol_b}
1/(1+t)
\pdfcraftend{slash_fraction@twocol_b}$ is often written with a slash.

\clearpage
\end{document}
